\section{Introduction}


Threshold signatures enable distributed signing authority where $t$ out of $n$ parties must collaborate to produce valid signatures, providing security against key compromise and single-point failures. While theoretical foundations exist, practical deployment across heterogeneous blockchain ecosystems presents significant engineering challenges:

\begin{enumerate}
\item \textbf{Protocol Diversity}: Different signature schemes (ECDSA, EdDSA, Schnorr) require distinct threshold protocols
\item \textbf{Chain Compatibility}: Each blockchain has unique signing requirements (message formats, encodings, hash functions)
\item \textbf{Dynamic Membership}: Real-world systems need to add/remove parties without reconstructing keys
\item \textbf{Performance}: Enterprise applications require sub-second signing latency
\item \textbf{Quantum Resistance}: Long-term security demands post-quantum alternatives
\end{enumerate}

This paper presents a unified framework addressing all challenges simultaneously through:

\paragraph{Universal Protocol Integration} Four complementary threshold protocols (CMP, FROST, LSS, Doerner) cover the complete spectrum from 2-party to large-scale threshold signing with ECDSA, EdDSA, and Schnorr support.

\paragraph{Multi-Chain Adapters} Native implementations for 20+ blockchains translate protocol outputs to chain-specific formats (BIP-340 Taproot, EIP-155/1559/4844 Ethereum, XRPL STX/SMT prefixes, Solana PDAs).

\paragraph{Dynamic Resharing} LSS protocol enables live membership changes and threshold updates without master key reconstruction or system downtime.

\paragraph{Post-Quantum Security} Corona lattice-based signatures provide quantum-resistant alternatives with Module-LWE hardness assumptions.

\paragraph{Production Validation} Comprehensive testing (100\% coverage, zero skipped tests), security audits, and real-world deployment securing billions in digital assets.

\subsection{Contributions}

\begin{enumerate}
\item \textbf{Unified Framework}: First production system integrating multiple threshold protocols with universal chain support
\item \textbf{LSS Dynamic Resharing}: Novel protocol for live membership changes with automatic fault tolerance
\item \textbf{Chain Adapter Architecture}: Extensible pattern for blockchain-agnostic threshold signatures
\item \textbf{Post-Quantum Integration}: Practical quantum-resistant threshold signatures for existing blockchains
\item \textbf{Performance Optimization}: Sub-25ms signing through constant-time arithmetic and parallel processing
\item \textbf{Production Deployment}: Battle-tested implementation securing enterprise custody and DeFi protocols
\end{enumerate}

\subsection{Version History}

\begin{itemize}
\item \textbf{v2021.02} (February 2021): Initial implementation with CMP and FROST protocols
\item \textbf{v2021.08} (August 2021): Added chain key derivation (BIP-32 compatible)
\item \textbf{v2023.05} (May 2023): Doerner 2-party protocol integration
\item \textbf{v2024.12} (December 2024): LSS dynamic resharing protocol
\item \textbf{v2025.08} (August 2025): Multi-chain adapters and Corona post-quantum signatures
\end{itemize}

